% !TeX root = ../main.tex
\section{Introduction}\label{sec:introduction}

% Why SLRs matter
\SLRs are an established research method in \SE.
Kitchenham and Charters~\cite{kitchenham2007guidelines} adapted the method from medicine, and later guidelines added systematic mapping studies and snowballing~\cite{petersen2015guidelines,wohlin2014guidelines}.
The core promise of an \SLR is that other researchers can follow each step.
Therefore, reviewers define a protocol up front, select papers with explicit criteria, and extract data with a fixed form.
Reporting guidelines such as \PRISMA~\cite{page2021prisma} then ask authors to disclose how many papers they found, screened, and excluded, and why.

% Problem
However, the data that an \SLR produces is hard to keep traceable.
Reviewers often record screening decisions and extracted values in spreadsheets, shared documents, or hosted services.
In such a setup, the link between an extracted value and the passage in the paper that supports it is often only in the reviewer's head.
Also, two or more reviewers usually extract the same papers independently, and someone has to compare and consolidate their answers.
With a shared spreadsheet, it is hard to see who changed which cell and why.
Hassler et al.~\cite{hassler2014outcomes} collected barriers to the \SLR process in a community workshop, and several of these barriers concern tool support.
Recently, \LLMs promise to take over parts of screening and extraction.
However, studies in \SE and other fields find that their output still needs human control~\cite{felizardo2024chatgpt,khraisha2024llm}.

% Gap
Several tools support parts of the process.
Mapping studies and surveys list many of them~\cite{marshall2013tools,marshall2014tools,harrison2020software}.
Many of these tools are hosted web applications that keep the review data on a server, for example, Covidence~\cite{covidence}, Rayyan~\cite{ouzzani2016rayyan}, and CADIMA~\cite{kohl2018cadima}.
Other tools focus on one activity, such as screening with active learning in ASReview~\cite{vandeschoot2021asreview}.
To the best of our knowledge, no tool combines local, plain-file review data with a version history, schema-driven extraction with evidence links, independent reviewers with a consolidation step, and \LLM proposals that a human must accept.

% Idea
In this paper, we present \SaiLoR, an open-source desktop tool that supports \SLRs in a systematic way.
The name contains the letters S, L, and R of \enquote{systematic literature review}.
The main idea is to treat review data like source code.
A review is a project of plain JSON files.
A schema describes what reviewers extract, and the tool validates every answer against it.
Each reviewer writes into their own files, and git records every change.
When reviewers pull each other's work, the tool merges their files field by field.
Optional \LLM assistance proposes answers with quotes from the paper, but the reviewer decides which proposal to apply, and the project records that an \LLM was involved.

% Running example
We use the example project that ships with \SaiLoR as a running example.
The project is a small review on \enquote{Code Search and Program Repair}.
A screening project holds eight candidate papers.
Five of them are included and carried into an extraction project with a nested schema.
\autoref{fig:annotate} shows this extraction project in the main view of \SaiLoR.

% Contributions
In summary, we make the following contributions:
\begin{itemize}
    \item We present \SaiLoR, a desktop tool that supports screening, data extraction, and consolidation of an \SLR with several reviewers.
    \item We describe the main architecture decisions behind \SaiLoR and their reasons, including local-first storage, git as collaboration layer, the move from Electron to Tauri, and the control of \LLM assistance.
    \item We report how we keep the tool traceable with 388 requirements that we linked to source code, tests, and documentation.
    \item We release \SaiLoR 2.0.0 as open source under the GPL-3.0 license on GitHub and Zenodo~\cite{sailor2026zenodo}.
\end{itemize}

% Outline
The remainder of the paper is structured as follows: In \autoref{sec:related-work}, we discuss related work.
In \autoref{sec:approach}, we explain which \SLR activities \SaiLoR supports and which design principles guide it.
\autoref{sec:tool} walks through the tool from the reviewer's view.
\autoref{sec:architecture} describes the main architecture decisions.
In \autoref{sec:requirements}, we present the requirements and their trace links.
We discuss limitations in \autoref{sec:discussion} and conclude in \autoref{sec:conclusion}.
